\documentclass[12pt]{article}
\usepackage[utf8]{inputenc}
\usepackage[T1]{fontenc}
\usepackage[margin=1in]{geometry}
\usepackage{setspace}
\usepackage{xcolor}
\usepackage[colorlinks=true, urlcolor=blue!60!black, linkcolor=blue!60!black]{hyperref}
\setlength{\parindent}{2em}
\setlength{\parskip}{1.25em}
\renewcommand{\baselinestretch}{1.0}
\begin{document}
\begin{center}
{\LARGE \textit{Wickard v. Filburn}, 317 U.S. 111 (1942)\par}
\vspace{0.8em}
{\large Case Brief\par}
\end{center}
\vspace{1em}

\section{Memory Jogger}

Under the Commerce Clause, Congress may regulate purely local, non-marketed economic activity---even a farmer's home-grown wheat consumed entirely on his own farm---when that activity, aggregated with the conduct of many others similarly situated, exerts a substantial economic effect on interstate commerce (\href{https://www.courtlistener.com/opinion/103716/wickard-v-filburn/}{CourtListener}; \href{https://supreme.justia.com/cases/federal/us/317/111/}{Justia}).

\section{Detailed Case Facts}

Roscoe C. Filburn owned and operated a small family farm in Montgomery County, Ohio, where he kept dairy cattle, sold milk, raised poultry, and sold poultry and eggs (\href{https://supreme.justia.com/cases/federal/us/317/111/}{Justia}; \href{https://www.oyez.org/cases/1940-1955/317us111}{Oyez}). Each fall he sowed a modest acreage of winter wheat for harvest the following July, customarily selling part of the harvest, feeding part to his livestock and poultry, grinding some into flour for his household, and retaining the rest as seed for the next planting.

To stabilize collapsing agricultural prices and prevent chronic surpluses during the Great Depression, Congress had enacted the Agricultural Adjustment Act of 1938 (``AAA''), 52 Stat. 31, 7 U.S.C. \S~1281 \textit{et seq.} Under the Act, the Secretary of Agriculture established annual national wheat acreage allotments, apportioned down to states, counties, and individual farms. Whenever projected total supply exceeded normal domestic consumption and exports by more than 35 percent, the Secretary proclaimed a compulsory national marketing quota, subject to approval by at least two-thirds of voting wheat farmers in a national referendum (\href{https://supreme.justia.com/cases/federal/us/317/111/}{Justia}).

In July 1940, federal agricultural officials notified Filburn that his 1941 wheat allotment was 11.1 acres with a normal yield of 20.1 bushels per acre. Ignoring the notice, Filburn sowed 23 acres in the fall of 1940---an excess of 11.9 acres. On May 19, 1941, shortly before the referendum, Secretary of Agriculture Claude R. Wickard delivered a radio address advocating approval of the quotas and highlighting higher federal wheat loans, without mentioning a pending statutory amendment enacted on May 26, 1941. That amendment increased the penalty on ``farm marketing excess'' from 15 cents per bushel to 49 cents per bushel (half the parity loan rate) and extended the penalty to excess wheat consumed on the farm. On May 31, 1941, wheat growers approved the quota in the referendum by an 81 percent to 19 percent margin (\href{https://supreme.justia.com/cases/federal/us/317/111/}{Justia}).

In July 1941, Filburn harvested and threshed his 23 acres, yielding 239 bushels of excess wheat from the 11.9 excess acres. At 49 cents per bushel, his penalty totaled \$117.11. Filburn refused to pay the penalty, store the excess under federal regulations, or turn it over to the Secretary; consequently, the county committee denied him a marketing card needed to sell his crop without exposing buyers to a lien. Filburn sued Secretary Wickard and local agricultural committee officials, arguing that Congress lacked Commerce Clause authority under Article I, \S~8, cl. 3 to regulate wheat grown solely for home consumption, and that the May 1941 amendment violated Fifth Amendment Due Process by retroactively increasing the penalty after he had already planted his crop (\href{https://www.courtlistener.com/opinion/103716/wickard-v-filburn/}{CourtListener}).

\section{Procedural History}

\begin{itemize}
    \item \textbf{Three-Judge U.S. District Court (S.D. Ohio):} In \textit{Filburn v. Helke}, 43 F. Supp. 1017 (S.D. Ohio 1942), a divided three-judge district court (2--1) ruled in Filburn's favor and permanently enjoined enforcement of the 49-cent penalty and crop lien (\href{https://supreme.justia.com/cases/federal/us/317/111/}{Justia}). The district court concluded that Secretary Wickard's radio speech invalidated the farmer referendum and that applying the May 1941 penalty increase to already-planted wheat was impermissibly retroactive under the Fifth Amendment (and inequitable).
    \item \textbf{U.S. Supreme Court:} Secretary Wickard and his codefendants took a direct appeal to the Supreme Court under 28 U.S.C. \S~380a. The case was initially argued on May 4, 1942, restored to the docket for reargument on October 13, 1942 to address the Commerce Clause question squarely, and decided on November 9, 1942 (\href{https://www.oyez.org/cases/1940-1955/317us111}{Oyez}; \href{https://supreme.justia.com/cases/federal/us/317/111/}{Justia}).
\end{itemize}

\section{Judicial Votes}

The decision was \textbf{unanimous (9--0)} (\href{https://supreme.justia.com/cases/federal/us/317/111/}{Justia}; \href{https://www.oyez.org/cases/1940-1955/317us111}{Oyez}). Justice Robert H. Jackson delivered the opinion of the Court, joined in full by Chief Justice Harlan Fiske Stone and Associate Justices Owen J. Roberts, Hugo L. Black, Stanley F. Reed, Felix Frankfurter, William O. Douglas, Frank Murphy, and James F. Byrnes. There were no concurring or dissenting opinions.

\section{Holding}

\textbf{Judgment:} Reversed. The Supreme Court vacated the district court's permanent injunction and upheld the wheat marketing quota and 49-cent-per-bushel penalty under both the Commerce Clause and the Fifth Amendment Due Process Clause (\href{https://supreme.justia.com/cases/federal/us/317/111/}{Justia}).

\textbf{Rule Announced:} Congress's power under the Commerce Clause reaches even purely local, intrastate activity that is not itself ``commerce''---including agricultural production intended solely for home consumption---if that activity exerts a substantial economic effect on interstate commerce, regardless of whether earlier doctrine would have labeled the effect ``direct'' or ``indirect.'' Furthermore, under the \textbf{aggregation principle}, an individual actor's trivial contribution does not exempt the activity from federal regulation when, taken together with that of many others similarly situated, the cumulative impact on interstate commerce is far from trivial (\href{https://www.courtlistener.com/opinion/103716/wickard-v-filburn/}{CourtListener}).

\textbf{Holding vs. Dicta:}
\begin{itemize}
    \item \textbf{Binding Holdings:} (1) The AAA's wheat marketing quotas and penalties, as applied to excess wheat consumed entirely on the farm, are valid exercises of the Commerce Clause because home-consumed wheat substantially affects interstate wheat supply and prices in the aggregate; (2) the May 26, 1941 amendment did not violate Fifth Amendment Due Process as retroactive legislation because the penalty attached only upon threshing after enactment and Filburn could have avoided it by feeding unthreshed wheat, storing the excess, or delivering it to the Secretary; and (3) Secretary Wickard's radio speech did not invalidate the farmer referendum absent evidence that voters were misled.
    \item \textbf{Dicta:} Justice Jackson's broad historical remarks invoking Chief Justice Marshall's suggestion in \textit{Gibbons v. Ogden} that effective restraints on the commerce power must come primarily from political rather than judicial processes went beyond what was needed to resolve the economic-effects issue, and later Supreme Court decisions---notably \textit{United States v. Lopez}, 514 U.S. 549 (1995), \href{https://supreme.justia.com/cases/federal/us/529/598/}{\textit{United States v. Morrison}}, 529 U.S. 598 (2000), and \href{https://supreme.justia.com/cases/federal/us/567/519/}{\textit{NFIB v. Sebelius}}, 567 U.S. 519 (2012)---reaffirmed judicially enforceable outer boundaries on Commerce Clause power.
\end{itemize}

\section{Analysis of Opinions}

\subsection{Unanimous Opinion of the Court (Justice Jackson)}

Justice Jackson's opinion dismantled the pre-1937 formalist barriers that had categorized ``production,'' ``manufacturing,'' and ``mining'' as strictly local matters whose effects on interstate commerce were deemed ``indirect'' as a matter of law (such as \textit{United States v. E. C. Knight Co.}, 156 U.S. 1 (1895), and \textit{Carter v. Carter Coal Co.}, 298 U.S. 238 (1936)). After swiftly reversing the district court's ruling on the Secretary's radio address as ``manifest error'' in Part I, Justice Jackson turned in Part II to the core Commerce Clause question (\href{https://supreme.justia.com/cases/federal/us/317/111/}{Justia}).

Tracing Commerce Clause doctrine from Chief Justice Marshall's expansive formulation in \textit{Gibbons v. Ogden} through Justice Holmes's pragmatic view in \textit{Swift \& Co. v. United States} and the \textit{Shreveport Rate Cases}, Justice Jackson replaced mechanical legal labels with a practical economic-effects inquiry. Whether Filburn's conduct was called ``production,'' ``consumption,'' or ``marketing'' was constitutionally immaterial; the decisive question was its real-world economic impact on the national wheat market.

Turning to the stipulated economic record, Justice Jackson explained that wheat was traded in a deeply integrated national and global market, where federal subsidies kept domestic farm prices around \$1.16 per bushel compared to a 40-cent world price. Crucially, home-consumed wheat represented the single most variable component in wheat disappearance, fluctuating by more than 20 percent of average production. Home-grown wheat affected interstate prices in two ways: first, marketable excess wheat ``overhangs the market'' and can flow into commercial channels if prices rise, checking price increases; second, even if never sold, home-grown wheat satisfies the grower's own demand for feed and flour, forestalling purchases on the open market and thereby competing with commercial wheat (\href{https://supreme.justia.com/cases/federal/us/317/111/}{Justia}).

Although Filburn's 239 bushels were minuscule in isolation, Justice Jackson articulated the landmark aggregation principle: an individual's trivial contribution cannot remove him from Congress's reach when the combined effect of thousands of similarly situated farmers is substantial. Conflicts between small diversified farmers and large commercial wheat growers were economic policy disputes reserved for Congress's legislative judgment, not judicial invalidation.

In Part III, Justice Jackson rejected Filburn's Fifth Amendment Due Process claims. Congress's regulatory scheme was neither arbitrary nor confiscatory, especially because cooperating farmers received substantial price-support benefits; as Jackson famously observed, it is hardly a denial of due process for the federal government to regulate what it subsidizes. Nor was the May 1941 amendment unconstitutionally retroactive: the penalty was triggered only when Filburn threshed his crop after enactment, and he had multiple penalty-free options, including storing the excess or feeding it unthreshed as hay (\href{https://www.courtlistener.com/opinion/103716/wickard-v-filburn/}{CourtListener}).

\section{Examples: Future Applications}

\begin{enumerate}
    \item \textbf{Same-Side (Valid Under \textit{Wickard}):} Congress prohibits the cultivation and possession of marijuana nationwide under a comprehensive controlled-substances regime. A patient in a state permitting medical cannabis grows six plants at home strictly for personal medicinal use and never buys or sells them. Following \textit{Wickard}, the Supreme Court upheld federal regulation in \href{https://supreme.justia.com/cases/federal/us/545/1/}{\textit{Gonzales v. Raich}}, 545 U.S. 1 (2005), because home-grown fungible commodities can rationally be deemed to affect interstate supply and demand in the aggregate.
    \item \textbf{Same-Side (Valid Under \textit{Wickard}):} Congress enacts a comprehensive federal statute capping carbon-intensive backyard coal smelting to stabilize national industrial metals markets and reduce interstate emissions. A small rural workshop smelts scrap iron purely for its own farm tools without ever selling a single piece. Because smelting is an economic production activity for a fungible commodity traded in an interstate market, aggregating all such self-producers exerts a substantial economic effect on interstate commerce under \textit{Wickard}.
    \item \textbf{Opposite-Side (Invalid---Beyond \textit{Wickard}):} Congress passes a federal criminal statute prohibiting the possession of a firearm within 1,000 feet of any local school, without any jurisdictional link to interstate commerce or connection to a broader economic regulatory scheme. As the Court held in \textit{United States v. Lopez}, 514 U.S. 549 (1995) (and reaffirmed for gender-motivated violence in \href{https://supreme.justia.com/cases/federal/us/529/598/}{\textit{United States v. Morrison}}, 529 U.S. 598 (2000)), \textit{Wickard}'s aggregation principle does not permit Congress to aggregate purely non-economic violent or possessory conduct to manufacture a general federal police power.
    \item \textbf{Opposite-Side (Invalid---Compelling Commerce):} Congress enacts a statute under the Commerce Clause requiring every adult resident who currently grows no vegetables to purchase at least \$200 of domestically grown produce each year to stabilize national agricultural prices. Under \href{https://supreme.justia.com/cases/federal/us/567/519/}{\textit{NFIB v. Sebelius}}, 567 U.S. 519 (2012), this exceeds \textit{Wickard}: Filburn was already actively engaged in the economic activity of farming wheat, whereas the Commerce Clause does not empower Congress to compel inactive individuals to enter commerce in the first place.
    \item \textbf{Fence-Sitter (Close Case):} Congress prohibits home distilling of spirits for purely personal, non-commercial consumption, defending the ban as necessary to protect federal liquor-tax revenues and regulate the interstate alcohol market. Challengers argue that unlike wheat in \textit{Wickard} or marijuana in \textit{Raich}, small-scale hobby distilling is non-commercial domestic activity historically outside interstate trade, and some federal courts have recently questioned Congress's enumerated power to ban home distilling (\href{https://alec.org/wp-content/uploads/2026/08/ALEC-AAF-LJC-26-93-Ream-Brief.pdf}{Litigation Overview}). Conversely, the government argues that home-distilled spirits substitute directly for commercially traded liquor in a heavily regulated national market, placing the case squarely on the modern boundary of \textit{Wickard} and \textit{Raich}.
\end{enumerate}

\section{Critique}

\textit{Wickard v. Filburn} is widely regarded as the high-water mark of federal Commerce Clause authority and remains a focal point of constitutional debate (\href{https://www.everycrsreport.com/reports/RL32844.html}{CRS Report RL32844}). From an \textbf{originalist and textualist perspective}, scholars such as Randy Barnett and Richard Epstein contend that \textit{Wickard} departed radically from the Founding-era public meaning of ``regulate Commerce . . . among the several States,'' which referred to trade and commercial intercourse across jurisdictional lines rather than purely local agriculture or home consumption (\href{https://fedsoc.org/commentary/videos/wickard-v-filburn-the-aggregation-principle-congressional-power-no-86}{Barnett, Federalist Society}; \href{https://fedsoc.org/fedsoc-review/the-meaning-of-regulate-commerce-to-the-constitution-s-ratifiers}{Federalist Society Review}). Critics further argue that pairing the ``substantial effects'' test with aggregation effectively collapses the doctrine of enumerated powers and erodes the Tenth Amendment, because virtually any human activity can be said to affect the national economy when aggregated across millions of people. Other conservative commentators note that \textit{Wickard} relied implicitly on an expansive reading of the Necessary and Proper Clause to reach non-commerce activities (\href{https://constitution.heritage.org/essays/a1-s8-c3-b/}{Heritage Guide to the Constitution}).

From a \textbf{progressive, pragmatic, and doctrinal perspective}, defenders view \textit{Wickard} as an essential constitutional adaptation to an integrated modern industrial economy. Scholars argue that the Framers designed Congress's enumerated powers to solve collective-action problems that individual states are incompetent to resolve on their own---precisely the interstate price-collapse and overproduction crisis addressed by the AAA (\href{https://www.scotusblog.com/2011/08/responding-to-the-eleventh-circuit-and-to-richard-epstein/}{SCOTUSblog}). Doctrinally, \textit{Wickard}'s rejection of artificial ``direct vs. indirect'' formalisms laid the groundwork for landmark federal legislation ranging from workplace protections to the Civil Rights Act of 1964, while later cases like \textit{Lopez}, \textit{Morrison}, and \textit{NFIB v. Sebelius} demonstrated that courts can preserve \textit{Wickard}'s economic-aggregation rule without granting Congress an unbounded police power (\href{https://southerncalifornialawreview.com/2019/12/08/the-limits-of-the-interstate-commerce-power-how-to-decide-the-close-cases-postscript-by-r-george-wright/}{Southern California Law Review}).

Underlying these critiques is a fundamental clash of \textbf{competing constitutional values}: limited federal government, state sovereignty, and formal textual boundaries on one side, versus national economic coherence, democratic deference to Congress on complex economic regulation, and functional realism on the other.

\section{Key Quotations}

\begin{quote}
``But even if appellee's activity be local, and though it may not be regarded as commerce, it may still, whatever its nature, be reached by Congress if it exerts a substantial economic effect on interstate commerce, and this irrespective of whether such effect is what might at some earlier time have been defined as `direct' or `indirect.''' \\
--- Justice Jackson, Opinion of the Court, 317 U.S. at 125 (\href{https://supreme.justia.com/cases/federal/us/317/111/}{Justia}; \href{https://www.courtlistener.com/opinion/103716/wickard-v-filburn/}{CourtListener})
\end{quote}

\begin{quote}
``That appellee's own contribution to the demand for wheat may be trivial by itself is not enough to remove him from the scope of federal regulation where, as here, his contribution, taken together with that of many others similarly situated, is far from trivial.'' \\
--- Justice Jackson, Opinion of the Court, 317 U.S. at 127--128 (\href{https://supreme.justia.com/cases/federal/us/317/111/}{Justia})
\end{quote}

\begin{quote}
``Home-grown wheat in this sense competes with wheat in commerce. The stimulation of commerce is a use of the regulatory function quite as definitely as prohibitions or restrictions thereon.'' \\
--- Justice Jackson, Opinion of the Court, 317 U.S. at 128 (\href{https://supreme.justia.com/cases/federal/us/317/111/}{Justia})
\end{quote}

\begin{quote}
``It is hardly lack of due process for the Government to regulate that which it subsidizes.'' \\
--- Justice Jackson, Opinion of the Court, 317 U.S. at 131 (\href{https://supreme.justia.com/cases/federal/us/317/111/}{Justia}; \href{https://www.courtlistener.com/opinion/103716/wickard-v-filburn/}{CourtListener})
\end{quote}

\end{document}
